\begin{tikzpicture}[
  >={Latex[length=3mm,width=2.2mm]},
  every node/.style={font=\small},
  uvec/.style={->,line width=1pt,black},
  mvec/.style={->,line width=1.3pt,black}
]

% Solenoide appiattito (forma a nastro chiuso irregolare)
% Contorno esterno
\draw[thick,black]
  plot [smooth cycle,tension=0.8] coordinates {
    (-4.5,0.8) (-3.0,2.2) (-0.5,2.0) (1.5,2.5) (3.5,1.8)
    (4.5,0.5) (3.5,-0.5) (2.0,0.2) (1.0,-0.2) (-0.5,-0.5)
    (-2.5,-0.8) (-4.0,-0.2)
  };
% Contorno interno (piccola distanza)
\draw[thick,black]
  plot [smooth cycle,tension=0.8] coordinates {
    (-4.0,0.7) (-2.7,1.8) (-0.5,1.6) (1.3,2.1) (3.0,1.5)
    (3.8,0.5) (3.0,-0.1) (2.0,0.5) (1.0,0.2) (-0.3,-0.1)
    (-2.2,-0.4) (-3.5,0.1)
  };

% Segmenti tra i due bordi (fili del solenoide)
\foreach \a in {0,20,...,340}{
  \draw[thin,black]
    ({3.5*cos(\a)+0.5},{1.2*sin(\a)+0.6}) --
    ({3.0*cos(\a)+0.5},{1.0*sin(\a)+0.6});
}

% Freccia corrente I
\draw[->,thick] (-3.5,2.3) -- ++(-0.8,0.3) node[above] {$I$};
\draw[->,thick] (-3.0,-1.2) -- ++(0.8,-0.3) node[below] {$I$};

% M e n-hat (a destra)
\draw[mvec] (4.2,0.3) -- ++(0,1.6) node[above] {$\mathbf{M}$};
\draw[uvec] (4.2,0.3) -- ++(1.0,0) node[right] {$\hat{n}$};

\end{tikzpicture}
